\documentclass[10pt,a4paper]{amsart}
\usepackage[margin=19mm]{geometry}
\usepackage{amsmath,amssymb,mathtools}
\usepackage[scr=rsfs,cal=euler]{mathalfa}
\usepackage{xfrac}
\usepackage[unicode,hidelinks,bookmarks=false]{hyperref}
\newcommand{\ie}{i.e.\ }
\newcommand{\cf}{cf.\ }
\newcommand{\resp}{respectively\ }
\newcommand{\Subsection}{\subsection}
\providecommand{\nobreakdash}{\nobreak\hbox{-}}
\setlength{\emergencystretch}{2em}
\multlinegap0pt
\usepackage{bm}
\usepackage{xcolor}
\newtheorem{lemma}{Lemma}
\newcommand{\hbeta}{\widehat\beta}
\newcommand{\lp}{\mathscr{L}}
\newcommand{\np}{\mathscr{N}}
\newcommand{\pp}{\mathscr{P}}
\newcommand{\rp}{\mathscr{R}}
\newcommand{\ts}{\mathcal{T}}
\title{Partizan Subtraction with full and truncated support}
\author{Ankita Dargad, Urban Larsson}
\date{Source: arXiv:2607.27989v1 (CC BY 4.0)}
\begin{document}
\maketitle

\noindent\emph{Source and licence.} Partizan Subtraction with full and truncated support. Version arXiv:2607.27989v1 (CC BY 4.0), distributed by arXiv under the Creative Commons Attribution 4.0 International licence, http://creativecommons.org/licenses/by/4.0/, as recorded in the arXiv licence metadata for this exact version and shown on the abstract page https://arxiv.org/abs/2607.27989v1. Full licence notice: \texttt{licenses/arxiv-2607.27989-CC-BY-4.0.txt}.\par\smallskip

\noindent\textit{Editorial scope.} Counted unit: the article's lemma block (source0.tex lines 3175--3194). The article's complete original proof of it is delivered with it (source0.tex lines 3195--3244, one \texttt{proof} environment of 50 lines). No mathematics is written, repaired or re-derived: the statement and the proof are the article's own lines, in the article's own notation, and no retained line is substituted or re-flowed.  Retained uncounted as the source-local context the proof needs: lines 1962--1969; lines 3131--3133.  Nothing was omitted from any retained span, so the package carries no omission record.  No retained line contains a citation, so the wrapper carries no bibliography and no bibliography entry.  No retained line contains a figure environment, an image-inclusion command or drawing code, so no image is delivered and the package declares no asset dependency.  Editorial formatting only, in addition: an AMS \texttt{amsart} preamble that restates the article's own declarations and macros used by the retained lines, namely the environments \texttt{lemma} on separate counters, together with the standard packages those lines need.  The retained environments print in this document as Lemma 1 in order of appearance, whereas the article prints the same items with the numbering of its own full document.  The complete native source of the article as distributed in the arXiv e-print for this exact version is bundled unchanged as \texttt{sources/206379/source0.tex}.
\begin{lemma}\label{lem: outcome TS heap<=b}
    Consider $G=\ts_\tau(n;a,b)$ where $a<b$ and $1\le \tau<b$. Then, 
    \begin{enumerate}
        \item\label{item:lem: outcome TS heap<=b:1} $o(G)=\lp$ if $n\le \tau$;
        \item\label{item:lem: outcome TS heap<=b:2} $o(G)=\np$ if $\tau+1\le n\le \min(\tau+a,b)$;
        \item\label{item:lem: outcome TS heap<=b:3} $o(G)=\rp$ if $\min(\tau+a,b)+1\le n\le b$.
    \end{enumerate}
\end{lemma}

 \begin{align}
    \hbeta=\left \lfloor \frac{b-a}{b-a-\tau}\right\rfloor.\label{eq:betahat}
\end{align} 

\begin{lemma}\label{lem: outcome ts below alpha block}
    Let $a,b$ and $\tau$ are positive integers such that $a<b$ and $\tau<b-a$. Then, for all $\beta<\hbeta$, %let 
    %     \begin{itemize}
    %     \item $L_\beta=(\beta-1)(b+1)+1$,
    %     \item $N_\beta=\beta(a+\tau+1)-a$, %=L_\beta-1+b-a-\beta(b-a-\tau)
    %     \item $R_\beta=\beta(a+\tau+1)$, %=N_\beta+a 
    %     \item $P_\beta=\beta(b+1)$. %=R_\beta+\beta(b-a-\tau)-1
    % \end{itemize}
    %Let %$\hbeta$ be such that $$\frac{(\hbeta-1)(b-a)-1}{\hbeta}<\tau\le \frac{\hbeta(b-a)-1}{\hbeta+1},$$ 
    %$\hbeta=\left \lfloor \frac{b-a}{b-a-\tau}\right\rfloor$
    %Then  %as in \eqref{eq:betahat},
    \begin{align*}
        o\big(\ts_\tau(n;a,b)\big)=\begin{cases}
            \lp & \text{ if } L_\beta \le  n<N_\beta ;\\
            \np & \text{ if } N_\beta\le  n< R_\beta;\\
            \rp & \text{ if } R_\beta\le  n< P_\beta;\\
            \pp & \text{ if } n=P_\beta.
        \end{cases}
    \end{align*}
\end{lemma}

\begin{proof}
    If $2\tau<b-a$, then $\hbeta=1$, and the statements holds vacuously. So, assume $2\tau\ge b-a$. We must prove that
    \begin{itemize}
        \item[(i)] Left wins $G$ if $L_\beta\le n< N_\beta$,
        \item[(ii)] First player wins $G$ if $N_\beta\le n< R_\beta$,
        \item[(iii)] Right wins $G$ if $R_\beta\le n< P_\beta$,
        \item[(iv)] Second player wins $G$ if $n=P_\beta$.
    \end{itemize}

    We first consider the case $\beta=1$. If $L_1\le n<P_1$, the statement holds using Lemma~\ref{lem: outcome TS heap<=b}. So let $n=P_1=b+1$: If Left starts, then any resulting position is either $\np$ or $\rp$ as $n-a>\tau+1=N_1$ and $n-1<P_1$. Thus, Right wins $G$. If Right starts, any resulting position is either $\lp$ or $\np$ as $n-b=1\ge L_1$ and $n-(\tau+1)\le a+\tau<R_1$. Thus, Left wins $G$. This completes the proof for $\beta=1$. 
    
    Now suppose $\beta>1$, i.e, $n\ge b+2$.

    % By induction, $L_{\beta-1}$, $N_{\beta-1},\;R_{\beta-1}$ and $P_{\beta-1}$ are as follows:
    % % Let $L_{\beta-1}$, $N_{\beta-1},\;R_{\beta-1}$ and $P_{\beta-1}$ are obtained from $L_\beta,N_\beta,R_\beta$ and $P_\beta$ respectively by replacing $\beta$ by $\beta-1$ in those. Thus
    % \begin{align*}
    %     L_{\beta-1}&=(\beta-2)(b+1)+1, \; & N_{\beta-1}&=(\beta-1)(a+\tau+1)-a,\\
    %     R_{\beta-1}&=(\beta-1)(a+\tau+1), \; & P_{\beta-1}&=(\beta-1)(b+1).
    % \end{align*}
    
    Consider item~(i).  Suppose first that Right starts in $G$ and moves to $G^R$. Denote the heap size of $G^R$ by $n_r$. Since $L_\beta\le n< N_\beta$, $$L_\beta-b\le n_r<N_\beta-(\tau+1).$$ By simple calculations, we get $L_\beta-b>L_{\beta-1}$ and $N_\beta-(\tau+1)=R_{\beta-1}$. Thus, $$L_{\beta-1}<n_r< R_{\beta-1}.$$ Hence, by induction, the outcome of $G^R$ is either $\lp$ or $\np$, and thus, Left wins playing first on $G^R$. Since $G^R$ is an arbitrary Right option, Left wins $G$. 
    
    Next suppose that Left starts in $G$. We claim that Left wins by removing 1 token. Let $G^L$ denote the Left option $\ts_\tau(n-1;a,b)$. If $n=L_\beta$, i.e., $n-1=P_{\beta-1}$, then outcome of $G^L$ is $\pp$ by induction and hence, Left wins $G$. Otherwise $L_\beta\le n-1< N_\beta-1$, and in this range, we already proved that Right loses as first player. Thus, Left wins $G$ by moving to $\ts_\tau(n-1;a,b)$.
    
    Consider item~(ii). First suppose that Left starts on $G$. Since the outcome of {\sc TS} is $\lp$ below the heap size $N_\beta$ by item~(i), Left wins if she can reach the heap size $N_\beta-1$. For this purpose, she needs to remove $n-N_\beta+1$ tokens %from the heap and wins as $o\left(\ts_\tau(N_\beta-1;a,b)\right)=\lp$ by item~(i). This is legal as 
    and this is possible as $$1\le n-N_\beta+1< R_\beta-N_\beta+1=a+1.$$ Therefore, Left wins $G$ by doing so. 
    
    Next suppose that Right starts on $G$. 
    %If $\beta=1$, then $\tau+1\le n\le \tau+a<b$, and Right wins by removing the full heap. Otherwise, 
    Right can win by reducing the heap size to $R_{\beta-1}$ by induction. For this purpose, Right needs to remove $n-R_{\beta-1}$ many tokens. Since $$\tau+1=N_\beta-R_{\beta-1}\le n-R_{\beta-1}<R_\beta-R_{\beta-1}=a+\tau+1\le b,$$ Right does so and wins $G$. % By induction, the outcome of $\ts_\tau(R_{\beta-1};a,b)$ is $\rp$, and hence Right wins.

    Consider item~(iii). Suppose Right starts on $G$. %If $\beta=1$, then $a+\tau+1\le n\le b$, and therefore, Right wins by removing the full heap. Otherwise, 
    By induction, Right wins if he can reach any heap size between $R_{\beta-1}$ and $P_{\beta-1}$, both included. We know \begin{align*}
        R_{\beta-1}&=R_\beta-(a+\tau+1),\\
        P_{\beta-1}&=P_\beta-(b+1).
    \end{align*}This implies $R_{\beta-1}+a+\tau+1\le n\le P_{\beta-1}+b$. Since all numbers from $a+\tau+1$ to $b$ are in the subtraction set of Right, he can always reach a heap size between $R_{\beta-1}$ and $P_{\beta-1}$, both included and thus, Right wins. 
    
    Now suppose Left starts on $G$ and her move leads to heap size $n_\ell$. Then $$N_\beta=R_\beta-a\le n_\ell< P_\beta-1.$$ We already proved that Right wins as first player on any of these heap sizes. Since, $n_\ell$ is arbitrary, Right wins $G$. 

    Consider item~(iv). Suppose Left starts on $G$ and move to heap size $n_\ell$. Then, $$N_\beta\le P_\beta-a\le n_\ell\le P_\beta-1.$$ We already proved that Right wins as first player on any of these heap sizes. Since $n_\ell$ is arbitrary, Right wins $G$. 
    
    Now suppose Right starts on $G$ and move to heap size $n_r$. Then, \begin{align}
        L_\beta\le P_\beta-b\le n_r&\le P_\beta-(\tau+1)\nonumber\\
        &=\beta(b+1)-(\tau+1)-R_\beta+R_\beta\nonumber\\
        &=\beta(b-a-\tau)-(\tau+1)+R_\beta\nonumber\\
        &<\frac{\tau+1}{b-a-\tau}(b-a-\tau)-(\tau+1)+R_\beta\label{eq:1}\\
        &=R_\beta\nonumber.
    \end{align}
    Equation~\eqref{eq:1} follows as $\beta\le \hbeta-1\le \frac{\tau}{b-a-\tau}<\frac{\tau+1}{b-a-\tau}$. We already proved that Left wins as first player on any of these heap sizes. Since, $n_r$ is arbitrary, Left wins $G$. 
\end{proof}

\end{document}
